% !TeX root = ../papers/nc-p3.tex
\section{特異な反対行列の大域Gram表示と二次錐上の八次鉛筆}
\label{sec:opposite-gram-pencil}
\begin{theorem}\label{thm:opposite-gram-pencil}
四辺有限な標準有限Frobenius窓の異なる反対根の行列 $N$ が一般階数3なら、
ある $4\times3$ 一次式行列 $U(x)$ により $N=UU^{\mathsf t}$ となる。
各 $0\ne t\in k^3$ に対して $M(t):x\mapsto U(x)t$ の階数は2以上である。
等方二次曲線 $T=\{t^{\mathsf t}t=0\}\simeq\mathbb P^1$ 上で
$\det M(t)$ は恒等零でない八次二項式となる。その三次小行列式全体の
$T$ 上の共通零因子の次数は $2,3,4$ のいずれかである。
さらに $E=X\cap\{\operatorname{rank} U\le1\}$ は長さ4以上であり、その最初の配置と
反対の配置における線型包は同じ次元（平面または三空間）を持つ。
\end{theorem}

\begin{proof}[一点の接触からの大域分解]
余因子を $ff^{\mathsf t}$ とし、前節で用いた $f$ の $X$ 上の評価が非零の一点を取る。
$p=(1,0,0,0)$、$w=(x_1,x_2,x_3)$ とし、出力基底で $N(p)$ の第一対角だけを1とする。
局所Schur補は二次から始まり、その二次部は $A B^{\mathsf t}+B A^{\mathsf t}$ である。
ここで $A,B$ は $w$ の一次式列で、単元係数の定数項は吸収した。
従って右下三行のブロック $H$ は大域的に
$LL^{\mathsf t}+AB^{\mathsf t}+BA^{\mathsf t}$ である。第一対角を
$x_0^2+\alpha x_0+c$、その下の列を $x_0L+b$ と書いた。
$\det H=f_0^2\ne0$ より $P=(L,A,B)$ は $k(w)$ で可逆で、
$H=P C P^{\mathsf t}$、$C$ は成分 $C_{11}=C_{23}=C_{32}=1$ の定数行列である。
$P^{-1}b=(a,r,s)^{\mathsf t}$ とすると $\det N=0$ より
$x_0^2+\alpha x_0+c=(x_0+a)^2+2rs$、特に $a=\alpha/2$ は一次式である。
$v=(x_0+a,L^{\mathsf t})^{\mathsf t}$ と置くと $N-vv^{\mathsf t}$ は一般階数2。
実際 $k(w)$ 上で残る列 $(s,A^{\mathsf t})^{\mathsf t}$ と
$(r,B^{\mathsf t})^{\mathsf t}$ の対称積である。
この大域的な多項式行列の二次形式は従って分裂し、一意分解環上のGaussの補題で
二つの多項式列の対称積とスカラー共通因子に書ける。どちらかの列が $x$ について次数零なら、
一般の像に固定方向が入り、$f$ の四成分が線型従属になる。斉次次数の和は2なので
両列とも一次式、共通因子は定数である。定数の非退化三変数形式を単位行列に合同変換して
所望の $U$ を得る。$f$ は符号を付けた四つの最大行列式の共通符号倍である。
$M(t)$ の階数が1なら $U(x)t$ は定方向の倍になり同じ従属性が生じる。
階数零も $U$ の一般階数3に反する。
\end{proof}

\begin{proof}[等方二方向の局所長]
$\Gamma\subset X\times T$ を $U(x)t=0$ で定める。各局所環 $D$ で一行のノルムが
単元となるよう出力基底を取れる。$N|_D$ は階数1で非零だからである。
その行に直交する平面の等方方向は二本で、$D$ 上で互いに異なる切断を与える。
実際この平面は非退化で、剰余体上の異なる二根をArtin局所環へ持ち上げればよい。
この二方向で残る三行を評価したものをそれぞれ $a_i,b_i$ とする。
二列 $a,b$ の対称積はSchur補の単元倍で
$a_i b_j+a_j b_i=0$ である。$I=(a_i),I'=(b_i)$ と置くと、
$\Gamma$ のこの二枚の局所長は $\alpha=\operatorname{length}(D/I)$,
$\beta=\operatorname{length}(D/I')$、$E$ の局所環は $F=D/(I+I')$ である。
また四つの最大行列式の生成するイデアル $J$ は $(a\times b)$ に等しく
$(I+I')J=0$。よって全体で $\operatorname{length}E\ge4$ となる。
$F$ が非零なら埋込次元は高々2である。実際被覆の二切断は周囲の完備正則局所環にも取れて、
そこで $I,I'$ を作る列がともに消える。両列の一次部までも零ならSchur補は三次以上になり、
前節までの三次接触排除に反する。

$D$ の完全跡 $\tau$ を取る。$I\times I'$ への積跡の制限の階数 $r$ は
$n-\alpha-\beta$ 以上である。左核は $I\cap\operatorname{ann} I'$ でその長さは高々 $\beta$ だからである。
一方 $I^2 I'=I I'^2=0$ なので、この対合を生成三列へ引き戻すと
$F^3\times F^3$ 上の交代対合となり、階数は同じである。
具体的には三つの線型汎関数 $\phi_i\in F^*$ による
$(u,v)\mapsto\sum_i\phi_i((u\times v)_i)$ である。
その階数は $2\operatorname{length}F$ 以下である。被約な点の和なら各点で交代三行の階数が2以下であり、
平面的な点は平坦にそのような点の和へ変形できる。有限自由な族で三つの汎関数も延長できるため、
小行列式の消滅が特殊繊維に移る。

ここで必要な平面上の平滑化を確認する。二変数局所環における有限余長イデアルは
Hilbert--Burch行列の最大小行列式で書ける。有限長なので多項式環の点の近傍で表示できる。
行列の各成分に独立な定数パラメータを加えれば、全空間から行列空間への射は滑らかである。
一段の階数降下は一般階数層で滑らかな余次元2、次の階数降下は余次元6である。
標数零の一般滑らかさと次元から、一般パラメータの曲面上の切断は滑らかな零次元となる。
元の孤立点を通る一般一変数変形について、完備局所環上の全切断は
期待余次元2でCohen--Macaulay、特殊繊維が有限なのでパラメータは非零因子である。
完備Nakayamaでその点の部分はパラメータ冪級数環上有限、従って自由で元の長さを保つ。
曲線状・長さ1の場合にも同じ評価が通る。
以上より
\[
 n\le \alpha+\beta+2\operatorname{length}F\le2(\alpha+\beta),
 \qquad 10\le\operatorname{length}\Gamma .
\]
\end{proof}

\begin{proof}[線型包と八次の零点]
$E$ 上で $U=pq^{\mathsf t}$ と書ける。$p$ は反対配置の射影座標で、
$q$ は射影二平面への射を与え、$q^{\mathsf t}q$ は単元である。
有限環上で直線束は自明化でき、階数1の行列の非零な成分でこの分解を取れるためである。
$T$ の一般の $t$ では $q^{\mathsf t}t$ が $E$ 上の単元なので、
$M(t)$ は最初の配置における $E$ の線型包を反対配置の線型包の上へ写す。
両配置で直線切断の長さは3以下なので、どちらも少なくとも平面である。
反対側の線型包が平面で方程式 $l$ を持てば、$l^{\mathsf t}U$ の三つの一次式は
最初の $E$ で消える。初めの線型包が三空間ならこれらは恒等零で、
$N$ の成分の独立性に反する。よって両側とも平面であり、
初めの平面の方程式 $m$ を用いて $l^{\mathsf t}U=m c^{\mathsf t}$、$c\ne0$ と書ける。
一般の $t\in T$ で $c^{\mathsf t}t\ne0$ なので、両平面の商の間でも $M(t)$ は同型である。
両線型包が三空間の場合も一般の $M(t)$ は既に可逆である。従って $T$ 上の行列式は非零で次数8である。

$T$ の行列式の零点の局所パラメータを $z$ とし、$M$ をその離散付値環でSmith形にする。
非零指数は高々二個で、$0\le a\le b$ と記す（階数3では $a=0$）。
行列式の消滅位数は $a+b$、三次小行列式の共通消滅位数は $a$。
$\Gamma$ の $z$ 冪フィルトレーションの $j$ 段目の商は $X$ 上の巡回層で、
Smith指数が $j$ 以下の全ての座標一次式に消される。
実際行列の局所基底で $z^{e_i}L_i(z,x)=0$ であり、$j\ge e_i$ なら
$z^j/z^{j+1}$ 上で $L_i(0,x)=0$ となる。これらの一次式は独立である。
従って $j<a$ の層は直線切断上で長さ3以下、$a\le j<b$ は一点上で長さ1以下、
以後は零である。その零点上の長さは $3a+(b-a)$ 以下。
$d$ を共通零因子の次数とすれば、全ての点の和から
$10\le\operatorname{length}\Gamma\le8+d$、
また $2d\le\sum(a+b)=8$ を得る。
\end{proof}
